\begin{figure*}[t] % t stands for top of the page, adjust as needed
    \centering
    \fbox{ % Use \fbox to draw a box around content
        \begin{minipage}{\textwidth} % Use minipage to limit the width of the box
            \centering
            \vspace{2pt}
            \textbf{\simulator{} can efficiently perform simulated gradient descent of an auxiliary model.} \\ % Adjust the title size and weight
            \vspace{4pt} % Add some vertical space between title and text
            \begin{theorem}
    Consider an auxiliary transformer with  $L$ layers, $\aux{D}$ embedding dimension, $\aux{H}$ attention heads, and a maximum sequence length of $\aux{T}$. 
    Given a hyperparameter $S$ (see~\Cref{sec:stack}), \simulator{} can perform an efficient forward pass~(\Cref{sec:exposition_linear_forward}), compute the simulated gradient~(\Cref{sec:modification}), and evaluate the updated auxiliary model with a total of $$\left(\frac{(c_1 S^2 + c_3) \aux{D}^2}{\min(\aux{H}, S^2)} \cdot \aux{D}^2 + c_2 S \aux{D} \min(S^2, \aux{H}) + c_3 \frac{\aux{T} \aux{D} S}{ \min(\aux{H}, S^2) } \right) L$$ parameters, with constants $c_1, c_2, c_3 < 150$.
    The \simulator{} model has $\simu{D} = S\aux{D}$ embedding dimension and $\simu{H}=\min(S^2, \aux{H})$  attention heads. 
    See~\Cref{tab:construction} for a detailed breakdown of the parameters.
    \label{thm:main}
\end{theorem} 
\vspace{2pt}
        \end{minipage}
    }
\end{figure*}